Inferring logical rules from bulk transcriptomics turns out to be difficult.
Logical rules perform poorly to predict PD-L1 gene expression across various attempted methods to infer them.
Qualitatively different rules perform similarly, making it difficult to infer a consistent set of rules.
The rules found for the training dataset, Broad, do not generalize well to the holdout dataset, Sanger.

As a basis, the method treated in~\cite{lawless_interpretable_2023} was used to
infer rulesets for PD-L1 expression from TF activity scores.
In this method, continuous data is first discretized via thresholding.
Then, an optimal ruleset is found via mathematical methods.
Here, ``optimal'' means achieving a high accuracy in predicting a target boolean.
The target boolean is whether PD-L1 expression exceeds a chosen threshold, representing whether PD-L1 is active or not.
A ruleset is a collection of ``rules'', each rule consisting of one or more ``threshold clauses''.
In our application, exceeding a threshold means that a TF exceeds a certain activity level.

While this method has been shown to work well on a variety of datasets, it does not work well on our dataset
(\autoref{ch:boolean-model}).
When allowing moderate ruleset complexity (equivalent to at most five three-TF rules, or at most six two-TF rules plus a one-TF rule),
accuracy is not high.
In some cases, the results prove that no meaningfully better ruleset of this form exists (\autoref{sec:boolean_model_appendix}).

Upon further inspection (\autoref{ch:threshold-model}),
simple rulesets perform only marginally worse than the more complex rulesets found before.
In particular, a rule with a single TF can yield a comparable accuracy.
This suggests that rulesets of this form are not a good predictor for PD-L1 expression.
Results also show that rulesets are not consistent when applying stratified cross-validation on the training dataset.

\clearpage

A novel variant of the base model, the ramp loss model (\autoref{ch:ramp_model}),
was designed to find more consistent rulesets and possibly discover new, previously hidden rules.
Both goals were not achieved; similar rulesets were found and
these were not consistent among splits of the training dataset during stratified cross-validation.

Additional experiments revealed that the base method works slightly better on tumors than on metastases.
An additional experiment also revealed that the method works slightly better on raw TF gene expression data,
rather than TF activity estimates, yielding qualitatively different results with slightly higher accuracy.
However, rulesets based on gene expression data are not expected to be biologically meaningful.
Finally, a TF exclusion experiment revealed that accuracy drops slightly upon exclusion of promising TF\@.

Overall, inferring logical rules predicting high PD-L1 expression from combinatorial TF activity
in the context of bulk transcriptomics was only marginally successful.
Recommendations for further research are discussed in \autoref{ch:recommendations}.